\section{微分可能なキーポイント--姿勢整合性}
\label{sec:consistency}

直接的なヒートマップ損失と姿勢損失は、異なるinterfaceを通じて同一のコートを教師する。
密予測キーポイントが姿勢予測の微分可能な射影と一致することを要求し、この共通の意味を
実行可能な形にする。

\begin{figure}[t]
  \centering
  \resizebox{\linewidth}{!}{\begin{tikzpicture}[
  font=\sffamily\footnotesize,
  box/.style={draw, rounded corners=2pt, align=center, minimum height=9mm,
    inner xsep=4pt, inner ysep=3pt, line width=0.55pt},
  densebox/.style={box, draw=gsblue, fill=gsblue!9},
  posebox/.style={box, draw=poseorange, fill=poseorange!9},
  lossbox/.style={box, draw=lossred, fill=lossred!8},
  auxbox/.style={box, draw=black!48, fill=softgray},
  arrow/.style={-{Latex[length=2.1mm,width=1.35mm]}, semithick,
    draw=black!72},
  densearrow/.style={-{Latex[length=2.1mm,width=1.35mm]}, semithick,
    draw=gsblue!88!black},
  posearrow/.style={-{Latex[length=2.1mm,width=1.35mm]}, semithick,
    draw=poseorange!90!black},
  maskarrow/.style={-{Latex[length=2.0mm,width=1.25mm]}, semithick,
    dashed, draw=black!52},
  lossarrow/.style={-{Latex[length=2.1mm,width=1.35mm]}, semithick,
    draw=lossred},
  grad/.style={-{Latex[length=2.2mm,width=1.4mm]}, semithick,
    dashed, draw=lossred},
  smalllabel/.style={font=\sffamily\scriptsize, align=center, text=black!72}
]
  \node[font=\sffamily\small\bfseries, text=gsblue!85!black, anchor=west]
    at (0,5.88) {Dense KP branch};
  \node[font=\sffamily\small\bfseries, text=poseorange!90!black, anchor=west]
    at (0,2.78) {Camera branch};

  \node[densebox, text width=15mm, minimum height=11mm] (logits)
    at (0.85,4.55)
    {KP logits $\mathbf z$\\[-0.15em]\scriptsize 14 maps\\[-0.2em]
     \scriptsize padded $H_{\max}\!\times\!W_{\max}$};
  \node[densebox, text width=3.0cm, minimum height=13mm] (softmax)
    at (3.40,4.55)
    {valid-region masked spatial softmax\\[-0.15em]
     \scriptsize $p_{bkyx}\propto\exp(z_{bkyx}/\tau)$\\[-0.15em]
     \scriptsize expectation $\sum_{y,x}p_{bkyx}[x,y]$};
  \node[densebox, minimum width=17mm, minimum height=10mm] (densecoord)
    at (6.25,4.55)
    {$\hat{\mathbf u}^{\mathrm{dense}}_{bk}$\\[-0.15em]
     \scriptsize continuous pixels};

  \node[auxbox, minimum width=25mm, minimum height=7mm] (validregion)
    at (3.40,6.22)
    {\scriptsize $V_b$: true image region\\[-0.2em]
     \scriptsize batch padding excluded};

  \draw[densearrow] (logits.east) -- (softmax.west);
  \draw[densearrow] (softmax.east) -- (densecoord.west);
  \draw[maskarrow] (validregion.south) -- (softmax.north);

  \node[posebox, text width=15mm, minimum height=11mm] (pose10d)
    at (0.85,1.57)
    {raw pose10D\\[-0.15em]\scriptsize query-head output};
  \node[posebox, text width=2.72cm, minimum height=14mm] (decode)
    at (3.30,1.57)
    {strict differentiable decode\\[-0.15em]
     \scriptsize centre $\hat{\mathbf C}$; row--6D $\xrightarrow{\rm GS}$
       $\hat{\mathbf R}\in\SOthree$\\[-0.15em]
     \scriptsize $\hat f=\exp(\widehat{\log f})>0$};
  \node[posebox, text width=3.15cm, minimum height=16mm] (project)
    at (6.57,1.57)
    {differentiable projection\\[-0.15em]
     \scriptsize $\mathbf X^{\rm cam}_{bk}=\hat{\mathbf R}^{\!\top}
       (\mathbf X_k-\hat{\mathbf C})$\\[-0.15em]
     \scriptsize $(\hat{\mathbf u}^{\rm pose}_{bk},\hat z_{bk})
       =\pi_{\hat f,\,\mathbf c^{\rm GT}}
       (\mathbf X^{\rm cam}_{bk})$};
  \node[auxbox, minimum width=30mm, minimum height=8mm] (known)
    at (6.57,0.10)
    {\scriptsize known canonical $\mathbf X_{1:14}$ in semantic order\\[-0.2em]
     \scriptsize augmentation-aware $\mathbf c^{\rm GT}=(c_x,c_y)$};

  \draw[posearrow] (pose10d.east) -- (decode.west);
  \draw[posearrow] (decode.east) -- (project.west);
  \draw[maskarrow] (known.north) -- (project.south);

  \node[lossbox, text width=4.55cm, minimum height=15mm] (coordloss)
    at (11.42,4.55)
    {\bfseries normalized coordinate consistency\\[-0.05em]
     \scriptsize $d_b=\sqrt{H_b^2+W_b^2},\quad
       r_{bk}=\lVert\hat{\mathbf u}^{\rm dense}_{bk}
       -\hat{\mathbf u}^{\rm pose}_{bk}\rVert_2/d_b$\\[-0.05em]
     \scriptsize $\displaystyle
       \mathcal L_{\rm coord}=
       \frac{\sum_{b,k}m^{\rm GT}_{bk}\,
       \operatorname{Huber}_{\delta}(r_{bk})}
       {\sum_{b,k}m^{\rm GT}_{bk}}$};
  \node[lossbox, text width=4.55cm, minimum height=14mm] (depthloss)
    at (11.42,1.57)
    {\bfseries cheirality term\\[-0.05em]
     \scriptsize $\displaystyle
       \mathcal L_{\rm depth}=
       \operatorname*{mean}_{m^{\rm GT}_{bk}=1}
       \operatorname{softplus}\!\left(
       \frac{z_{\min}-\hat z_{bk}}{s_z}\right)$\\[-0.05em]
     \scriptsize invalid predicted depth is penalized, never masked};
  \node[auxbox, text width=4.35cm, minimum height=8mm] (visibility)
    at (11.42,6.22)
    {\scriptsize fixed $m^{\rm GT}_{bk}$: generated supervision / in-frame authority\\[-0.2em]
     \scriptsize independent of predicted score, coordinate, and depth};

  \draw[densearrow] (densecoord.east) -- (coordloss.west);
  \draw[posearrow] (project.east) -- (depthloss.west);
  \draw[posearrow] (project.east) -- (9.05,1.57) |-
    ($(coordloss.south west)+(0.35,0)$);
  \draw[maskarrow] (visibility.south) -- (coordloss.north);
  \draw[maskarrow] (visibility.east) -- (13.92,6.22) |-
    (depthloss.east);

  \node[lossbox, text width=2.30cm, minimum height=16mm] (total)
    at (15.56,3.16)
    {\bfseries auxiliary objective\\[-0.05em]
     \scriptsize $\mathcal L_{\rm kp\leftrightarrow pose}$\\[-0.1em]
     \scriptsize $=\mathcal L_{\rm coord}
       +\lambda_z\mathcal L_{\rm depth}$\\[0.15em]
     \scriptsize standard: both-grad};
  \draw[lossarrow] (coordloss.east) -- (total.north west);
  \draw[lossarrow] (depthloss.east) -- (total.south west);
  \node[smalllabel, text width=2.6cm, text=lossred!90!black] at (15.56,2.02)
    {direct KP and pose losses remain};

  \draw[grad, rounded corners=3pt] (total.north) -- (15.56,7.02) --
    node[smalllabel, above, pos=0.58, text=lossred!90!black]
      {$\nabla_{\mathbf z}\mathcal L$: gradient to dense branch}
    (0.85,7.02) -- (logits.north);
  \draw[grad, rounded corners=3pt] (total.south) -- (15.56,-0.58) --
    node[smalllabel, below, pos=0.58, text=lossred!90!black]
      {$\nabla_{\rm pose10D}\mathcal L$: gradient to camera branch}
    (0.85,-0.58) -- (pose10d.south);
\end{tikzpicture}
}
  \caption{\textbf{分岐間の幾何学的目的関数。}  密予測分岐は有効領域上のspatial softmaxと
  期待値を用いる一方、姿勢分岐は予測された中心、回転、および焦点距離により、物理的順序を
  持つ3次元コート点を射影する。生成時の教師適格性により、どの残差を数えるかを固定する。
  正規化Huber項が画像座標を整合させ、cheirality項がカメラ後方に配置された点を罰する。
  デフォルトの目的関数は両分岐へ勾配を伝播する。}
  \label{fig:consistency}
\end{figure}

\subsection{paddingを考慮した連続座標}

batch collationは、各sampleをbatch内の最大の高さと幅までpaddingする。このcanvas全体に対する
spatial softmaxは正しくない。有効logitが負の場合、定数0のpadding logitでさえ確率質量を
吸収し得るためである。\(z_{bkyx}\) をbatch要素 \(b\)、物理キーポイントチャネル \(k\)、画素
\((x,y)\) に対するlogitとし、sampleの実際の拡張後画像に属するBoolean領域を \(V_b\) とする。次式を
計算する。
\begin{align}
  p_{bkyx}
    &=\frac{\mathbb{1}[(y,x)\in V_b]\exp(z_{bkyx}/\tau)}
      {\sum_{(y',x')\in V_b}\exp(z_{bky'x'}/\tau)}, \\
  \widehat{\mathbf{u}}^{\mathrm{dense}}_{bk}
    &=\sum_{(y,x)\in V_b}p_{bkyx}[x,y]^{\top},
  \label{eq:masked_expectation}
\end{align}
ここで \(\tau>0\) である。この経路にはsigmoid、top-\(k\)、rounding、argmax、CPU transfer、
またはdetachは存在しない。maskはBooleanであり、logitとdeviceおよび空間shapeを共有し、
すべてのsampleについて少なくとも1画素の有効画素を含まなければならない。不正な入力は
reduction前に例外を発生させる。

この期待値は、target-court singleton \KP contractに限定して適用する。multi-courtかつmulti-peakな
チャネルに1つの期待値を適用すると、物理的意味を持たないコート間の座標が生成される。

\subsection{予測カメラからの射影}

\(\mathbf{X}_{bk}\) を、姿勢教師に格納された厳密な物理的順序における既知の正準3次元座標と
する。\cref{eq:pose10d,eq:rotation_decode}で復号された姿勢は、カメラ中心
\(\widehat{\mathbf{C}}_b\)、canonical-from-camera回転 \(\widehat{\mathbf{R}}_b\)、および焦点距離
\(\widehat f_b\) を与える。ground-truthの並進、回転、焦点距離を用いずに変換および射影する。
\begin{align}
  \widehat{\mathbf{X}}^{\Cam}_{bk}
    &=[\widehat X_{bk},\widehat Y_{bk},\widehat Z_{bk}]^{\top}
      =\widehat{\mathbf{R}}_b^{\top}
      (\mathbf{X}_{bk}-\widehat{\mathbf{C}}_b), \\
  \widehat{\Kmat}_b
    &=\begin{bmatrix}
       \widehat f_b&0&c_{x,b}\\
       0&\widehat f_b&c_{y,b}\\
       0&0&1
      \end{bmatrix}, \\
  \overline Z_{bk}
    &=\operatorname{sgn}_{+}(\widehat Z_{bk})
      \max(|\widehat Z_{bk}|,z_{\min}), \\
  \widehat{\mathbf{u}}^{\mathrm{pose}}_{bk}
    &=\pi_{z_{\min}}(\widehat{\Kmat}_b,
       \widehat{\mathbf{X}}^{\Cam}_{bk})
      =\begin{bmatrix}
        \widehat f_b\widehat X_{bk}/\overline Z_{bk}+c_{x,b}\\
        \widehat f_b\widehat Y_{bk}/\overline Z_{bk}+c_{y,b}
       \end{bmatrix},
  \label{eq:predicted_projection}
\end{align}
ここで、\(\operatorname{sgn}_{+}(z)=1\) は \(z\geq0\) のとき、そうでないときは \(-1\) である。
設定された \(z_{\min}>0\) はsafe-depthの大きさであると同時にcheirality thresholdでもある。
姿勢ヘッドは主点を予測しないため、主点 \((c_{x,b},c_{y,b})\) はaugmentation-awareなtarget
intrinsicsから取得する。これは幾何metadataであり、漏洩した姿勢教師信号ではない。
有界な座標射影は不正なdepthでも有限値を維持し、後述のcheirality項がその補正勾配を与える。

射影の前に、adapterは姿勢教師内の \texttt{semantic\_to\_physical} と、すべての密予測チャネルの
singleton physical indexが厳密に等しいことを検証する。また、両分岐が同一の拡張後model-pixel
座標系を用いることも検証する。nearest-neighbour matchingや暗黙のsource-to-model rescalingは
許容しない。

\subsection{教師信号で固定されたロバスト損失}

\(m^{\mathrm{GT}}_{bk}\) は、生成されたin-front、in-frame、およびrenderer-supportの正とする情報に
基づき、target-court pointが教師対象として適格であることを表す。重要な点として、このmaskは
予測キーポイントスコアおよび予測depthから独立している。画像対角長
\(d_b=\sqrt{H_b^2+W_b^2}\) を用いると、正規化残差と座標損失は
\begin{align}
  r_{bk}
    &=\frac{\|\widehat{\mathbf{u}}^{\mathrm{dense}}_{bk}
              -\widehat{\mathbf{u}}^{\mathrm{pose}}_{bk}\|_2}{d_b}, \\
  \mathcal{L}_{\mathrm{coord}}
    &=\frac{\sum_{b,k}m^{\mathrm{GT}}_{bk}
      \operatorname{Huber}(r_{bk};\delta)}
      {\sum_{b,k}m^{\mathrm{GT}}_{bk}}.
  \label{eq:coordinate_loss}
\end{align}
である。スカラーHuber関数の定義として
\(
  \operatorname{Huber}(r;\delta)=\tfrac12r^2
\)
を \(r\leq\delta\) の場合に用い、それ以外の場合は
\(
  \delta(r-\tfrac12\delta)
\)
とする。適格な点を持たないsampleは座標残差に寄与しない。ただし、batch全体で適格な点が
0個の場合は、静かにゼロ損失とするのではなくconfiguration/data errorとする。

適格な点をカメラ後方へ配置しても、教師信号を無効化できてはならない。そこで、次式を加える。
\begin{equation}
  \mathcal{L}_{\mathrm{cheir}}=
  \frac{\sum_{b,k}m^{\mathrm{GT}}_{bk}
    \operatorname{softplus}\!\left(
      \frac{z_{\min}-\widehat Z_{bk}}{s_z}\right)}
       {\sum_{b,k}m^{\mathrm{GT}}_{bk}},
  \qquad
  \mathcal{L}_{\mathrm{aux}}=
     \mathcal{L}_{\mathrm{coord}}+w_z\mathcal{L}_{\mathrm{cheir}},
  \label{eq:cheirality}
\end{equation}
ここで \(z_{\min},s_z>0\) である。学習の全目的関数は
\begin{equation}
  \mathcal{L}_{\mathrm{total}}=
  \mathcal{L}_{\mathrm{dense}}+\mathcal{L}_{\mathrm{pose}}
  +w_{\mathrm{aux}}(t)\mathcal{L}_{\mathrm{aux}}.
  \label{eq:total_loss}
\end{equation}
である。設定可能なweight 0のwarmup後、\(w_{\mathrm{aux}}(t)\) を宣言された値まで線形に増加
させる。これにより、未学習の2つの分岐が互いに唯一のteacherとなることを防ぎながら、
学習全体を通して直接的な教師信号を維持する。標準modeでは両分岐を通じて勾配を伝播する。
2つの明示的なablationでは、それぞれ姿勢側または密予測側の勾配を停止するが、forward residualは
変更しない。

\subsection{局所交点shortcutの制約}
\label{sec:shortcut}

\begin{figure}[t]
  \centering
  \resizebox{\linewidth}{!}{\begin{tikzpicture}[
  font=\sffamily\footnotesize,
  arrow/.style={-{Latex[length=2.1mm,width=1.35mm]}, semithick,
    draw=black!70},
  posearrow/.style={-{Latex[length=2.1mm,width=1.35mm]}, semithick,
    draw=poseorange!90!black},
  densearrow/.style={-{Latex[length=2.1mm,width=1.35mm]}, semithick,
    draw=gsblue!88!black},
  callout/.style={-{Latex[length=1.8mm,width=1.15mm]}, thin,
    draw=lossred!85!black},
  card/.style={draw=black!28, fill=white, rounded corners=2pt,
    line width=0.5pt},
  note/.style={draw=lossred, fill=lossred!7, rounded corners=2pt,
    align=center, inner xsep=5pt, inner ysep=3pt},
  smalllabel/.style={font=\sffamily\scriptsize, align=center, text=black!72}
]
  \draw[draw=black!16, fill=softgray!45, rounded corners=3pt]
    (0,0) rectangle (7.18,5.86);
  \draw[draw=black!16, fill=softgray!45, rounded corners=3pt]
    (7.48,0) rectangle (16.90,5.86);

  \node[font=\sffamily\small\bfseries, text=lossred!88!black, anchor=west]
    at (0.28,5.51) {(a) Local appearance is ambiguous};
  \node[font=\sffamily\small\bfseries, text=courtgreen!85!black, anchor=west]
    at (7.77,5.51) {(b) Global render-supervised training pressure};

  \draw[card] (0.35,1.13) rectangle (3.23,4.98);
  \draw[card] (3.94,1.13) rectangle (6.82,4.98);
  \node[font=\sffamily\scriptsize\bfseries] at (1.79,4.67) {view / court A};
  \node[font=\sffamily\scriptsize\bfseries] at (5.38,4.67) {view / court B};

  \fill[courtgreen!88]
    (0.67,2.55) -- (2.93,2.55) -- (2.57,4.28) -- (1.03,4.28) -- cycle;
  \draw[white, line width=0.5pt]
    (0.67,2.55) -- (2.93,2.55) -- (2.57,4.28) -- (1.03,4.28) -- cycle;
  \draw[white, line width=0.4pt]
    (1.02,2.55) -- (1.27,4.28)
    (2.58,2.55) -- (2.33,4.28)
    (0.82,3.23) -- (2.79,3.23)
    (1.15,3.75) -- (2.45,3.75)
    (1.80,3.23) -- (1.80,3.75);

  \fill[courtgreen!88]
    (4.23,2.49) -- (6.54,2.69) -- (6.04,4.31) -- (4.53,4.16) -- cycle;
  \draw[white, line width=0.5pt]
    (4.23,2.49) -- (6.54,2.69) -- (6.04,4.31) -- (4.53,4.16) -- cycle;
  \draw[white, line width=0.4pt]
    (4.58,2.52) -- (4.77,4.18)
    (6.18,2.66) -- (5.80,4.29)
    (4.36,3.17) -- (6.35,3.34)
    (4.65,3.69) -- (6.17,3.83)
    (5.36,3.26) -- (5.40,3.76);

  \fill[lossred] (1.13,3.23) circle[radius=1.5pt];
  \draw[lossred!80!black, line width=0.45pt] (1.13,3.23) circle[radius=3.2pt];
  \fill[lossred] (5.90,3.30) circle[radius=1.5pt];
  \draw[lossred!80!black, line width=0.45pt] (5.90,3.30) circle[radius=3.2pt];

  \draw[fill=courtgreen!92, draw=lossred!80!black, rounded corners=1pt]
    (1.24,1.55) rectangle (2.34,2.35);
  \draw[white, line width=0.75pt]
    (1.79,1.62) -- (1.79,2.28)
    (1.31,2.04) -- (2.27,2.04);
  \draw[fill=courtgreen!92, draw=lossred!80!black, rounded corners=1pt]
    (4.83,1.55) rectangle (5.93,2.35);
  \draw[white, line width=0.75pt]
    (5.38,1.62) -- (5.38,2.28)
    (4.90,2.04) -- (5.86,2.04);
  \draw[callout] (1.13,3.17) -- (1.67,2.36);
  \draw[callout] (5.90,3.24) -- (5.50,2.36);
  \node[smalllabel, text=lossred!85!black] at (1.79,1.34)
    {CourtKP14 role $k$};
  \node[smalllabel, text=lossred!85!black] at (5.38,1.34)
    {CourtKP14 role $k'\ne k$};
  \node[smalllabel, font=\sffamily\scriptsize\bfseries] at (3.59,0.76)
    {same local ``T'' pattern};
  \node[smalllabel] at (3.59,0.40)
    {$\phi(\mathcal N_A)\!\approx\!\phi(\mathcal N_B)$ locally,
     but the global semantic roles differ};

  \draw[card] (7.79,1.48) rectangle (10.27,4.94);
  \node[smalllabel, font=\sffamily\scriptsize\bfseries]
    at (9.03,4.64) {accepted render poses};
  \fill[courtgreen!78, draw=courtgreen!85!black]
    (8.45,2.38) rectangle (9.61,4.07);
  \draw[white, line width=0.45pt]
    (8.54,2.48) rectangle (9.52,3.97)
    (8.54,3.22) -- (9.52,3.22)
    (9.03,2.48) -- (9.03,3.97);
  \node[regular polygon, regular polygon sides=3, rotate=18,
    draw=poseorange!90!black, fill=poseorange!28, minimum size=4.0mm,
    inner sep=0pt] at (8.18,3.70) {};
  \node[regular polygon, regular polygon sides=3, rotate=145,
    draw=poseorange!90!black, fill=poseorange!28, minimum size=4.0mm,
    inner sep=0pt] at (9.89,4.24) {};
  \node[regular polygon, regular polygon sides=3, rotate=205,
    draw=poseorange!90!black, fill=poseorange!28, minimum size=4.0mm,
    inner sep=0pt] at (10.00,2.28) {};
  \node[regular polygon, regular polygon sides=3, rotate=-18,
    draw=poseorange!90!black, fill=poseorange!28, minimum size=4.0mm,
    inner sep=0pt] at (8.15,2.15) {};
  \draw[black!30, thin] (8.25,3.64) -- (9.03,3.25);
  \draw[black!30, thin] (9.82,4.14) -- (9.18,3.47);
  \draw[black!30, thin] (9.92,2.36) -- (9.16,2.96);
  \draw[black!30, thin] (8.23,2.24) -- (8.91,2.95);
  \node[smalllabel] at (9.03,1.79)
    {broad camera and\\court-instance coverage};

  \node[draw=poseorange, fill=poseorange!9, rounded corners=2pt,
    minimum width=19mm, minimum height=22mm, align=center] (global)
    at (11.48,3.18) {};
  \node[font=\sffamily\scriptsize\bfseries, text=poseorange!90!black]
    at (11.48,4.00) {global task encoder};
  \foreach \x/\y in {-0.47/0.25,-0.16/0.25,0.15/0.25,0.46/0.25,
                       -0.47/-0.05,-0.16/-0.05,0.15/-0.05,0.46/-0.05} {
    \node[draw=gsblue, fill=gsblue!18, rounded corners=0.7pt,
      minimum width=2.4mm, minimum height=2.4mm, inner sep=0pt]
      at ($(global.center)+(\x,\y)$) {};
  }
  \node[circle, draw=poseorange!90!black, fill=poseorange!30,
    minimum size=5.3mm, inner sep=0pt, font=\sffamily\scriptsize\bfseries]
    at ($(global.center)+(0,-0.52)$) {$q$};
  \node[smalllabel] at (11.48,1.82)
    {whole-frame context\\$+$ global query};

  \draw[card] (12.68,1.42) rectangle (16.63,4.98);
  \node[smalllabel, font=\sffamily\scriptsize\bfseries,
    text=courtgreen!85!black] at (14.73,4.68)
    {one global projected\\CourtKP14 configuration};
  \fill[courtgreen!87]
    (13.19,2.22) -- (16.37,2.22) -- (15.91,4.22) -- (13.69,4.22) -- cycle;
  \draw[white, line width=0.48pt]
    (13.19,2.22) -- (16.37,2.22) -- (15.91,4.22) -- (13.69,4.22) -- cycle;
  \draw[white, line width=0.4pt]
    (13.52,2.22) -- (13.89,4.22)
    (16.04,2.22) -- (15.71,4.22)
    (13.36,2.83) -- (16.23,2.83)
    (13.56,3.61) -- (16.06,3.61)
    (14.80,2.83) -- (14.81,3.61);

  \foreach \x/\y in {
      13.19/2.22,16.37/2.22,13.69/4.22,15.91/4.22,
      13.52/2.22,16.04/2.22,13.89/4.22,15.71/4.22,
      13.36/2.83,16.23/2.83,14.80/2.83,
      13.56/3.61,16.06/3.61,14.81/3.61} {
    \draw[gsblue, line width=0.55pt] (\x,\y) circle[radius=2.05pt];
    \fill[poseorange] (\x,\y) circle[radius=1.15pt];
  }
  \node[smalllabel] at (14.73,1.78)
    {\textcolor{poseorange}{$\bullet$} pose projection \quad
     \textcolor{gsblue}{$\circ$} dense KP expectation};

  \draw[arrow] (10.27,3.18) -- (global.west);
  \draw[posearrow] ($(global.east)+(0,0.28)$) -- (12.68,3.46);
  \draw[densearrow] ($(global.east)+(0,-0.28)$) -- (12.68,2.90);

  \node[note, text width=8.15cm] at (12.27,0.63)
    {\scriptsize Across diverse renders, globally inconsistent local outputs
     are penalized by pose and eligible-point consistency targets.\\[-0.05em]
     \scriptsize\itshape This is training pressure, not proof that local
     representations are eliminated.};
\end{tikzpicture}
}
  \caption{\textbf{局所交点shortcutの制約。}  白線交差の周囲の小さな近傍は、複数の物理的identityと
  姿勢に整合し得る。レンダリングされた姿勢の多様性は、そのような局所的観測情報を保持しながら
  全体配置を変化させる。直接的な姿勢教師信号、大域クエリアテンション、物理channel identity、
  および1つの共有された投影 \KP 配置は、互いに矛盾する局所出力を罰する。これらは局所的な内部表現を
  不可能にするものではなく、実ドメインにおけるshortcutの減少は依然として経験的な問題である。}
  \label{fig:shortcut}
\end{figure}

おおむね直交する2本の白線を含むcropを考える。その局所appearanceは、service-line intersection、
baseline/sideline corner、手前のコートではなく奥のコート、あるいは異なる射影ビューのいずれとも
整合し得る。純粋に局所的なclassifierは、学習集合内の頻度biasを利用して、これらのidentityから1つを
選択できる。

\Method{}は、相互に接続された3つの段階でこのshortcutに対処する。
第1に、3DGSレンダリングは、正確な物理identityと6DoFを保持しながら、azimuth、range、height、
画角、target court、partial coverage、およびbackground contextを変化させる。その結果、
同じ局所motifが異なる大域配置内に現れ得る。第2に、位置を付与しないcourt queryは
空間patch集合全体へattentionでき、大域姿勢を回帰するよう直接学習される。2D RoPEは、その観測情報が
現れた位置を保持する。第3に、\cref{eq:coordinate_loss}は、同一の投影剛体コートと一致しない、
教師対象として適格なすべての密予測チャネルを罰する。適格なjunctionを1つ独立に移動させれば残差が増加する。
一方、姿勢予測と密予測を同時に移動させても、\cref{eq:direct_losses}の直接損失によって
引き続き罰せられる。

この機構は局所特徴を禁止しない---lineとjunctionは引き続き有用な画像情報である---うえ、密予測分岐が
どの内部特徴を用いるかを決定することもできない。その代わり、これらの出力が、大域的に整合し直接教師された
カメラ--コートの説明を支持できない場合には損失を増加させる。したがって、依存度が実際に減少するかを検証する
には、\cref{sec:experiments}のshortcutに焦点を当てた診断が必要である。正確な姿勢ラベルによって、この出力
levelの制約を低コストで実現できる。同じjoint supervisionを人手で作るには、意図的に中心から外したビューや
部分被覆ビューを含め、すべての学習画像を較正する必要がある。
